\begin{lemma}[A polynomial restricted to a circle is a trigonometric polynomial]
Let $n,m\in\mathbb{N}$, let $P\in\mathbb{R}[X_1,\dots,X_n]$ have total degree
$\deg P\le m$, and let $A,B\in\mathbb{R}^n$. Then the function
\[
  t\ \longmapsto\ P\bigl(\cos(t)A+\sin(t)B\bigr)
\]
is a real trigonometric polynomial of degree at most $m$.
\end{lemma}
\begin{proof}
Write $\gamma(t)=\cos(t)A+\sin(t)B$, so that coordinatewise
\begin{equation}\label{eq:gctp-coord}
  \gamma(t)_a=\cos(t)A_a+\sin(t)B_a\qquad(1\le a\le n),
\end{equation}
and let $V_N=\{T:\mathbb{R}\to\mathbb{R}\ \text{a real trigonometric polynomial
of degree at most }N\}$ in the sense of \noderef{TrigPolyLe}.

\emph{Step 1: $V_N$ is closed under passing to a larger degree.} Let $k\le N$ and
$T\in V_k$, with witnesses $\alpha,\beta:\mathbb{N}\to\mathbb{R}$, so
$T(t)=\sum_{j=0}^{k}(\alpha_j\cos(jt)+\beta_j\sin(jt))$ for all $t$. Define
$\alpha'_j=\alpha_j$ and $\beta'_j=\beta_j$ for $j\le k$, and
$\alpha'_j=\beta'_j=0$ for $j>k$. Then for every $t$
\[
  \sum_{j=0}^{N}\bigl(\alpha'_j\cos(jt)+\beta'_j\sin(jt)\bigr)
  =\sum_{j=0}^{k}\bigl(\alpha_j\cos(jt)+\beta_j\sin(jt)\bigr)=T(t),
\]
because every term with $k<j\le N$ vanishes. So $(\alpha',\beta')$ witnesses
$T\in V_N$. (This step is not vacuous: the witnesses supplied by $T\in V_k$ need
not themselves vanish above $k$, which is why they are replaced.)

\emph{Step 2: each coordinate function lies in $V_1$.} Fix $a$. Taking
$\alpha_0=0$, $\alpha_1=A_a$, $\beta_0=0$, $\beta_1=B_a$ and all further entries
$0$, we get, for every $t$,
\[
  \sum_{j=0}^{1}\bigl(\alpha_j\cos(jt)+\beta_j\sin(jt)\bigr)
  =\alpha_0+\alpha_1\cos t+\beta_1\sin t
  =\cos(t)A_a+\sin(t)B_a=\gamma(t)_a ,
\]
using $\cos(0\cdot t)=1$ and $\sin(0\cdot t)=0$. So $t\mapsto\gamma(t)_a$ lies in
$V_1$ by \noderef{TrigPolyLe}.

\emph{Step 3: a monomial of degree $k$ gives an element of $V_k$.} Let
$d\in\mathbb{N}^n$ with $k=\sum_a d_a$. We claim
$t\mapsto\prod_{a=1}^{n}\gamma(t)_a^{\,d_a}$ lies in $V_k$. The product is a
product of exactly $k$ factors, each of them one of the functions
$t\mapsto\gamma(t)_a$ (the factor $t\mapsto\gamma(t)_a$ repeated $d_a$ times).
We argue by induction on $k$. If $k=0$ the product is the constant function $1$,
which lies in $V_0$: take $\alpha_0=1$ and all other entries $0$, so that
$\sum_{j=0}^{0}(\alpha_j\cos(jt)+\beta_j\sin(jt))=\alpha_0=1$. If $k\ge1$, split
off one factor: the product equals $g\cdot h$ where $g(t)=\gamma(t)_{a_0}$ for
some index $a_0$ with $d_{a_0}\ge1$, and $h$ is the product of the remaining
$k-1$ factors. By Step 2, $g\in V_1$; by the induction hypothesis, $h\in V_{k-1}$.
By \noderef{TrigPolyMul} the pointwise product $g\cdot h$ lies in
$V_{1+(k-1)}=V_k$, which completes the induction.

\emph{Step 4: each monomial of $P$ gives an element of $V_m$.} Let
$d\in\operatorname{supp}P$. By the background proposition on the monomial
expansion, $\sum_a d_a\le\deg P\le m$. So by Step 3 the function
$t\mapsto\prod_a\gamma(t)_a^{\,d_a}$ lies in $V_k$ with $k=\sum_ad_a\le m$, and
hence in $V_m$ by Step 1.

\emph{Step 5: assembling the finite sum.} By the background proposition on the
monomial expansion, for every $t$
\begin{equation}\label{eq:gctp-exp}
  P\bigl(\gamma(t)\bigr)
  =\sum_{d\in\operatorname{supp}P}\bigl(\operatorname{coeff}_dP\bigr)
    \prod_{a=1}^{n}\gamma(t)_a^{\,d_a}.
\end{equation}
We show by induction on the cardinality of a finite set
$S\subseteq\operatorname{supp}P$ that
$t\mapsto\sum_{d\in S}(\operatorname{coeff}_dP)\prod_a\gamma(t)_a^{d_a}$ lies in
$V_m$. For $S=\emptyset$ the function is the constant $0$, which lies in $V_m$
(take all witness entries $0$). For the inductive step write
$S=S_0\cup\{d_0\}$ with $d_0\notin S_0$; then the function for $S$ is
$1\cdot G_0+(\operatorname{coeff}_{d_0}P)\cdot G_1$ where
$G_0(t)=\sum_{d\in S_0}(\operatorname{coeff}_dP)\prod_a\gamma(t)_a^{d_a}$ lies in
$V_m$ by the induction hypothesis and
$G_1(t)=\prod_a\gamma(t)_a^{(d_0)_a}$ lies in $V_m$ by Step 4. By
\noderef{TrigPolyLinComb}, applied to $G_0$ and $G_1$ with the scalars
$c=1$ and $d=\operatorname{coeff}_{d_0}P$, that linear combination lies in
$V_m$; and it agrees pointwise with the function for $S$, hence equals it.

Taking $S=\operatorname{supp}P$ and using \eqref{eq:gctp-exp}, the function
$t\mapsto P(\gamma(t))$ agrees at every $t$ with a function in $V_m$, hence
equals it and so lies in $V_m$. That is the assertion.
\end{proof}
